\begin{tikzpicture}[font=\scriptsize, >=latex]
% array vs pointer codegen
\draw[thick, fill=accentsoft] (0.3,2.8) rectangle (4.6,4.6);
\node[font=\footnotesize] at (2.45,4.2) {\ttfamily int a[4]};
\node[align=left, font=\tiny, text width=3.9cm] at (2.45,3.4)
  {\texttt{a[i]} $\to$ base $+$ i$\times$4,\\
   one load from a fixed offset ---\\
   the array \emph{is} its address};
\draw[thick, fill=graybg] (5.5,2.8) rectangle (9.8,4.6);
\node[font=\footnotesize] at (7.65,4.2) {\ttfamily int *p};
\node[align=left, font=\tiny, text width=3.9cm] at (7.65,3.4)
  {\texttt{p[i]} $\to$ load p first,\\
   then base $+$ i$\times$4 ---\\
   two memory trips where the array needed one};
\node[align=left, text=slate, text width=10cm] at (5.05,2.0)
  {``Arrays are pointers'' is a half-truth: the array's address is
   compile-time, the pointer's is runtime --- the generated code keeps
   the receipt.};
\end{tikzpicture}
